\begin{mainquestion}{csdesign}{Common Source Amplifier Design}

A CS amplifier must be designed to provide an output signal with 1.0 V amplitude using a load $R_{D} = 82\ \text{k}\Omega$. Note: This is a design problem; please provide lower \& upper bounds for answers to parts (b)-(d).

\begin{partquestion}{mingm}{If a voltage gain $\left|A_{v}\right| \geq 14\ \frac{\text{V}}{\text{V}}$ is desired, what minimum transconductance $g_{m}$ is required? You may use this $g_{m}$ value to answer the parts below.}
\begin{answer}
If $\lambda = 0$, the gain magnitude of a common source amplifier is:

\[
\left|A_{v}\right| = g_{m} R_{D} \geq 14
\]

so

\[
g_{m} \geq \frac{14}{R_{D}} = \frac{14}{82} = 0.171\ \text{mS}.
\]
\end{answer}
\end{partquestion}

\begin{partquestion}{idanvov}{Using a DC supply $V_{DD} = 3.3\ \text{V}$, what values of $I_{D}$ and $V_{OV} = V_{GS} - V_{t}$ would you choose?}
\begin{answer}
\begin{kaitlyn}

\noindent\fbox{%
\begin{minipage}{\dimexpr\textwidth-2\fboxsep-2\fboxrule\relax}
        \textbf{What this question is asking:} Design a bias point where the resulting output voltage's AC component has enough headroom to swing upwards without hitting the DC rail and  swing downwards without leaving the saturation regime.
\end{minipage}%
}%

Recall our saturation condition:

\[\boxed{V_{DS} \geq V_{OV}}\]

At the DC operating point, the drain voltage is

\[V_{DS,Q} = V_{DD} - I_D R_D.\]

However, $V_{DS}$ changes when an AC input signal is applied. Since the common-source amplifier has negative gain,

\[A_v = -g_mR_D,\]

the worst case for remaining in saturation occurs when the input is at its maximum positive excursion $(+v_{gs})$. At this instant:

\begin{itemize}
    \item the output swings downward, so the drain source voltage becomes $V_{DS,Q}-v_{ds}$, and
    \item the gate-source voltage increases, so the instantaneous overdrive becomes $V_{OV}+v_{gs}$.

\end{itemize}

Therefore, at the worst-case point in the signal swing, we require

\[\boxed{V_{DD} - I_D R_D - v_{ds}\geq V_{OV} + v_{gs} }\]

Notice that this equation contains quantities from both the DC bias point and the AC small-signal model. We now move between these levels of abstraction to rewrite the constraint in terms of $V_{OV}$.

From the MOSFET bias relationship,

\[ g_m = \frac{2I_D}{V_{OV}} \qquad\Rightarrow\qquad I_D = \frac{1}{2}g_mV_{OV}.\]

From the amplifier gain,

\[ |A_v| = \frac{|v_{ds}|}{|v_{gs}|} \qquad\Rightarrow\qquad v_{gs} = \frac{v_{ds}}{|A_v|}.\]

Substituting these relationships into the saturation constraint gives
\end{kaitlyn}

\begin{align*}
V_{DD} - \frac{1}{2}g_{m}V_{OV}R_{D} - v_{ds} &\geq V_{OV} + \frac{v_{ds}}{A_{v}} \\
V_{DD} - v_{ds} - \frac{v_{ds}}{A_{v}} &\geq V_{OV}\left(1 + \frac{1}{2}g_{m}R_{D}\right) \\
V_{OV} &\leq \frac{V_{DD} - v_{ds} - \frac{v_{ds}}{A_{v}}}{1 + \frac{1}{2}g_{m}R_{D}}\\
V_{OV} &\leq \frac{3.3 - 1.0 - \frac{1.0}{14}}{1 + \frac{1}{2}(0.171)(82)}\\
V_{OV} &\leq 0.279\ \text{V}
\end{align*}

and

\begin{align*}
    I_D &= \frac{1}{2} g_m V_{OV} \leq \frac{1}{2} g_m (0.279) \\
    &\leq \frac{1}{2} (0.171)(0.279) \\
    &\leq \qty{0.0238}{\milli\ampere} \\
\end{align*}

At the high end, we can't be too close to the rail in order to allow the small signal output to vary across its maximum amplitude. $V_{DS} \leq V_{DD} - v_{ds} = 3.3 - 1.0 = 2.3\ \text{V}$. The constraints on the current becomes:

\[
I_{D} \geq \frac{V_{DD} - V_{DS}}{R_{D}} = \frac{3.3 - 2.3}{82} = 0.0122\ \text{mA}
\]

and

\[
V_{OV} = \frac{2I_{D}}{g_{m}} \geq \frac{2(0.0122)}{0.171} = 0.143\ \text{V}.
\]

Putting all this together gives the acceptable current range of $12\ \mu\text{A} \leq I_{D} \leq 24\ \mu\text{A}$ and acceptable overdrive voltage range of $0.143\ \text{V} \leq V_{OV} \leq 0.279\ \text{V}$.
\end{answer}
\end{partquestion}

\begin{partquestion}{wlratio}{What $\frac{W}{L}$ ratio is required if $k^{\prime} = \mu_{n} C_{ox} = 70\ \frac{\mu\text{A}}{V^{2}}$?}
\begin{answer}
Since $g_{m} = k'\frac{W}{L}V_{OV}$,

\[
\frac{W}{L} = \frac{g_{m}}{k^{\prime} V_{OV}} = \frac{0.171}{(0.070)V_{OV}} = \frac{2.44}{V_{OV}}
\]

Given the range of $V_{OV}$ from part (\cref{prob:csdesign:idanvov}), we can see the corresponding range for $\frac{W}{L}$: $8.76 \leq \frac{W}{L} \leq 17.08$.
\end{answer}
\end{partquestion}

\begin{partquestion}{vgsrange}{If $V_{t} = 0.6\ \text{V}$, find the range of $V_{GS}$.}
\begin{answer}
Since $V_{GS} = V_{OV} + V_{t} = V_{OV} + 0.6$, we can use the range of $V_{OV}$ from part (\cref{prob:csdesign:idanvov}) to find $0.743\ \text{V} \leq V_{GS} \leq 0.879\ \text{V}$.
\end{answer}
\end{partquestion}

\end{mainquestion}
